\section{海岸不是陆上征服的边缘}

厦门与金门之间的水道并不因为陆上府县易手而安静下来。顺治十年（1653年），郑成功以两岛为支点，把军队、商船和仍可调动的沿海关系拢在同一套经营中。岛屿提供停泊与转运的地点，港口把粮食、人员和消息送向舰队；但要把这些东西送到前线，还得越过潮汐、水道与清军据守的城镇。郑氏因此不能只守海面：漳州、泉州、同安一线的攻守，关系到腹地粮食能否抵达海边，也关系到城防能否把这条补给线截断。

顺治十一年至十三年间，围攻漳州以及在泉州、同安附近争夺水陆通道，正是这种压力的展开。郑军需要扩大粮饷与招募的基础，清军则依托府县城防和陆路援兵，使其难以迅速取得大陆府城。对港口周边的人们而言，战争不是两支舰队在远海相遇：征粮、避乱、治安的变化，会落在村落、渡口、盐场和市集不同的位置。我们可以从军报、地方志与郑氏文书看见攻守和若干地点，却很难据此排出每一处居民如何选择、何时离开或继续交易的整齐次序。商人、船民、耕作者、守城者和逃亡者面对的风险不必相同，更不能被概括成一种一致的“支持”或“抵抗”。

这也是郑氏海上力量的两面。商船和侨民网络能够输送筹饷的机会，岛屿也能让军政机构避开陆上包围；可是船只、兵员和收入的数字多见于敌对军报或后出的郑氏叙事，不能把它们当作精确的家底。较为稳妥的判断是：福建海疆已无法仅由陆上州县来控制。清廷后来把海禁、迁界与水师建设连在一起，并不是出于抽象的禁海观念，而是要处理港口、岛屿、腹地与消息网络彼此相连的难题。

顺治十六年（1659年）北上南京，使这种难题沿江而上。郑成功试图利用长江下游的粮赋与明遗民网络，把海上基地转化为更大的反清战场。围城并不只考验冲锋；军队要在内河环境中维持粮食、船只、城外营地和联络，清军则集中江防、城门与守备力量。郑军未能突破城防而撤回，显示其攻势在内河攻坚与长期补给面前难以持续。清廷记录、杨英《从征实录》与荷兰海上报告都能确认北伐、围城和撤退的主干，但内应的范围、作战计划及伤亡互有歧异；把这次行动写成几乎成功的复明，或写成毫无准备的冒险，都会超出材料所能承担的程度。

撤退没有使台湾自动成为唯一的去处，却使那里成为更迫切的战略选择。顺治十八年（1661年），清廷命福建、广东沿海迁界，意在把郑氏同岸上的粮食、人口与信息来源隔开。谕令所能直接告诉我们的，是中央选择了以强制迁徙制造隔离带；它不能替我们证明每一县迁了多远、多少人死于途中，或秘密贸易有多大规模。沿海县志和通志留下的迁撤、复界线索，使渔盐、耕作、宗族与港口关系被打断的代价可以进入叙事；但这些记载的密度和视角各异，不能把局部遭遇拼成一幅均一的海岸图景。

迁界也并未把海面与陆地彻底分开。禁令增加了往来成本，海岸居民、地方官和海上经营者仍可能在不同地点以不同方式应对：有人被迫内迁，有人守着可耕之地或亲族关系，有人继续依赖港口和航路。现有材料足以说明隔离带未能完全阻绝往来，却不足以量化这些变通，更不足以把它们叙述成全海岸一致、持续不断的走私网络。政策的直接军事目的，是压缩郑氏补给的可能；它的副作用，则是把本已被战争牵动的生计、迁徙和地方关系再度推入不确定之中。

同年四月，郑成功军登陆台湾，围攻荷兰东印度公司在热兰遮城的据点。对郑氏而言，这并非单纯向外扩张：南京行动受挫之后，需要一处能够安置军政、取得粮食并控制航路的基地；对岛上的荷兰据点、汉人移民社会与原住民社群而言，战争则把既有的土地、征发与交往关系重新置于强制之下。围城的胜负取决于守军补给能否进入、周边村社与航道由谁控制，以及攻守双方能否把人力和粮食留在城下，而不只是城墙前的一次决战。

次年正月，热兰遮城守军投降，荷兰东印度公司在台湾西南的统治至此终结。投降条约和《热兰遮城日志》能够确定交接的时点，也能让人看见围城中荷兰方面的处境；清廷军报与后出的史书则呈现北京获知这一消息的路径。它们不能替城外居民、被俘者或原住民社群作证。兵员、联盟和伤亡在中荷战报中各有立场，因而不能用任一方的数字填满这场转换。可以确认的是，郑氏接收了城堡、田地和贸易节点，台湾由此成为反清政权筹集人口与财政资源的基地；至于这种接收怎样抵达岛上不同社群、各处又怎样回应，仍须依赖保存不均的地方、契约和多语材料逐地重建。

郑成功在康熙元年五月去世，使新得的基地立刻面对继承问题。城堡已经易手，并不表示军政、宗族与海商网络已经形成稳定的服从关系；郑经最终掌权，是在这些网络重新安排权力与资源的过程中出现的结果。有关死因、继承争执和军中处置的细节，主要依赖后出的郑氏材料与荷兰报告，叙事只能把握重组的压力，不能替当事人补出一致的动机。

康熙二年（1663年），清军联合部分郑氏降将攻取厦门、金门，郑氏主力退守台湾。熟悉海上条件的降将与沿海据点帮助清军夺取近大陆的前进基地，但厦金易手并不等于清廷已经取得制海权。对郑氏而言，失去两岛迫使台湾的土地、港口和跨海贸易承担更重的军政负担；对清廷而言，迁界所制造的隔离、沿海守备与水师建设仍要持续付出社会与后勤成本。这里的制度得失并不在于哪一方终于“控制了海岸”，而在于双方都不得不以迁徙、城防、航路和粮饷来维持一场尚未结束的海上战争。
